\section{Cariotipo}

Il \textbf{cariotipo} è la rappresentazione ordinata dei cromosomi metafasici di un individuo: le cellule vengono bloccate in metafase, colorate (ad esempio con Giemsa) e fotografate al microscopio; i cromosomi vengono poi ritagliati dalla foto e appaiati con i relativi omologhi.
\begin{itemize}
  \item i cromosomi da 1 a 20 vengono ordinati in base alla grandezza decrescente; il cromosoma 21 precede il 22, anche se il cromosoma 21 è più piccolo;
  \item con tecniche di colorazione a fluorescenza (cariotipo spettrale), ogni coppia di cromosomi viene marcata con una diversa sonda fluorescente, per facilitarne l'identificazione e l'appaiamento secondo l'ordine convenzionale.
\end{itemize}

% ---------------------------------------------------------------------------
\subsection{Ereditarietà mendeliana nell'uomo}

\rubrica{Eredità autosomica dominante}
\begin{itemize}
  \item il carattere si manifesta sia in omozigosi che in eterozigosi;
  \item il rischio di ricorrenza è del 50\% (ciascun figlio/a di un individuo affetto ha il 50\% di probabilità di essere affetto);
  \item ogni individuo affetto ha un genitore affetto;
  \item individui non affetti non trasmettono il carattere ai figli.
\end{itemize}

\rubrica{Eredità autosomica recessiva}
\begin{itemize}
  \item il carattere si manifesta solo in omozigosi;
  \item entrambi i genitori dell'individuo che manifesta il carattere sono portatori dell'allele (eterozigoti);
  \item nel caso di un carattere associato a malattia genica, il rischio di ricorrenza per ciascun figlio/a di portatori sani è di 1 su 4 (25\%);
  \item tanto più è rara la malattia, tanto più è probabile una consanguineità nei genitori.
\end{itemize}

\rubrica{Alberi genealogici}
Rappresentano l'ereditarietà mendeliana negli individui umani, con simboli convenzionali (quadrato = maschio, cerchio = femmina, simbolo colorato = individuo con il carattere; numeri romani = generazioni, numeri arabi = individui di una generazione).
\begin{itemize}
  \item \textbf{carattere dominante}: compare in ogni generazione;
  \item \textbf{carattere recessivo}: gli individui affetti sono spesso figli di consanguinei.
\end{itemize}

% ---------------------------------------------------------------------------
\subsection{I cromosomi sessuali}

\begin{itemize}
  \item le femmine hanno due cromosomi X e producono gameti (uova), ciascuno con un cromosoma X;
  \item i maschi hanno un cromosoma X e uno Y, e producono gameti (spermatozoi) metà dei quali portano un cromosoma X, l'altra metà un cromosoma Y;
  \item i maschi trasmettono il loro cromosoma Y ai figli, non alle figlie; i figli maschi ereditano il loro cromosoma X dalla madre.
\end{itemize}

\rubrica{Eredità legata al sesso}
Esempio storico: l'eredità dell'emofilia nei discendenti della Regina Vittoria d'Inghilterra, portatrice del gene mutato recessivo, trasmesso ai discendenti maschi.

\rubrica{Inattivazione del cromosoma X}
L'inattivazione di un cromosoma X compensa gli effetti della doppia dose di geni presente nelle femmine dei mammiferi.
\begin{itemize}
  \item nello zigote e nelle divisioni precoci dell'embrione, entrambi i cromosomi X sono attivi;
  \item durante lo sviluppo cellulare, a caso un cromosoma X si inattiva in ciascuna cellula, creando una linea cellulare specifica;
  \item una volta che un certo cromosoma X si è inattivato, rimane tale in tutte le cellule figlie discendenti $\rightarrow$ la femmina è un mosaico per i geni X-linked;
  \item esempio: il fenotipo ``calico'' (a chiazze arancioni e nere) nelle gatte eterozigoti per gli alleli del colore del pelo arancione e nero.
\end{itemize}

% ---------------------------------------------------------------------------
\subsection{Alterazioni cromosomiche}

Alterazioni della struttura dei cromosomi.

\rubrica{Delezione}
Perdita di un segmento del cromosoma. Può causare gravi problemi se il segmento cromosomico perso contiene geni essenziali per il normale sviluppo o per le normali funzioni cellulari.

\rubrica{Duplicazione}
Un segmento si stacca da un cromosoma e si inserisce sul suo omologo, nel quale gli alleli inseriti si sommano a quelli già presenti. Può avere effetti benefici o nocivi a seconda dei geni e degli alleli localizzati nel segmento duplicato. Esempio: l'emoglobina, a seguito di duplicazioni cromosomiche e successive mutazioni dei geni duplicati, ha dato origine a nuove varianti utili.

\rubrica{Traslocazione reciproca}
Un segmento si stacca da un cromosoma e si inserisce su un altro cromosoma non omologo. Esempio: il linfoma di Burkitt, in cui geni che controllano la divisione cellulare, nella nuova collocazione, si trovano vicino a regioni che controllano geni molto attivi $\rightarrow$ maggiore espressione $\rightarrow$ divisioni cellulari incontrollate $\rightarrow$ sviluppo del cancro.

\rubrica{Inversione}
Un segmento si stacca e si reinserisce sullo stesso cromosoma, ma ruotato di 180°, in modo tale che l'ordine dei geni è invertito. Alcuni geni possono subire rotture interne con conseguente perdita della loro specifica funzione.

% ---------------------------------------------------------------------------
\subsection{Mutazioni genomiche: non-disgiunzione meiotica}

La non-disgiunzione è il mancato distacco dei cromosomi omologhi (o dei cromatidi fratelli) durante la meiosi.
\begin{itemize}
  \item \textbf{non-disgiunzione nella prima divisione meiotica}: determina la migrazione allo stesso polo del fuso di entrambi i cromosomi omologhi (ciascuno bicromatidico) di una coppia $\rightarrow$ si ottengono due gameti con un cromosoma in più ($n+1$) e due gameti con un cromosoma in meno ($n-1$);
  \item \textbf{non-disgiunzione nella seconda divisione meiotica}: produce due gameti normali ($n$), uno con un cromosoma in più ($n+1$) e uno con un cromosoma in meno ($n-1$).
\end{itemize}

\rubrica{Aneuploidia}
Condizione irregolare del numero di cromosomi, in cui la variazione è limitata a uno o a pochi elementi dell'assetto cromosomico, per eccesso o per difetto.

\rubrica{Trisomia del cromosoma 21 (sindrome di Down)}
Presenza di tre copie del cromosoma 21. L'incidenza della sindrome di Down aumenta con l'età della madre.

\rubrica{Non-disgiunzione di cromosomi sessuali}
Combinazioni anomale di cromosomi sessuali prodotte da non-disgiunzione dei cromosomi X nella femmina:

\begin{table}[htbp]
  \centering
  \small
  \renewcommand{\arraystretch}{1.3}
  \begin{tabularx}{\textwidth}{|c|c|X|}
    \hline
    \textbf{Combinazione} & \textbf{Frequenza approssimativa} & \textbf{Effetti} \\
    \hline
    XO & 1 su 5.000 nati & sindrome di Turner; femmine con ovaie iposviluppate, sterilità, intelligenza e genitali esterni normali, bassa statura e seno poco sviluppato \\
    \hline
    XXY & 1 su 2.000 nati & sindrome di Klinefelter; maschi con testicoli iposviluppati, sterili, intelligenza di solito normale, peli radi, mammelle più sviluppate del normale; caratteristiche simili sono presenti negli individui XXXY e XXXXY \\
    \hline
    XYY & 1 su 1.000 nati & sindrome XYY; maschi apparentemente normali, spesso più alti della media \\
    \hline
    XXX & 1 su 1.000 nati & sindrome del triplo X; femmine apparentemente normali, con funzioni mentali normali o solo lievemente ritardate \\
    \hline
  \end{tabularx}
\end{table}

\rubrica{Condizioni ereditarie nell'uomo}

\begin{table}[htbp]
  \centering
  \small
  \renewcommand{\arraystretch}{1.2}
  \begin{tabularx}{\textwidth}{|X|X|}
    \hline
    \multicolumn{2}{|l|}{\textbf{Eredità autosomica recessiva}} \\
    \hline
    Albinismo & assenza di pigmentazione (melanina) \\
    \hline
    Lobi dell'orecchio attaccati & nessuno \\
    \hline
    Fibrosi cistica & eccesso di muco nei polmoni e nelle cavità digerenti \\
    \hline
    Anemia falciforme & danni a diversi tessuti e organi \\
    \hline
    Galattosemia & danni a cervello, fegato e occhi \\
    \hline
    Fenilchetonuria & ritardo mentale \\
    \hline
    Malattia di Tay--Sachs & ritardo mentale, morte \\
    \hline
    \multicolumn{2}{|l|}{\textbf{Eredità legata all'X}} \\
    \hline
    Emofilia A & insufficienza nella coagulazione del sangue \\
    \hline
    Daltonismo & incapacità di distinguere il rosso dal verde \\
    \hline
    Sindrome della femminilizzazione testicolare & assenza di organi sessuali maschili, sterilità \\
    \hline
    \multicolumn{2}{|l|}{\textbf{Cambiamenti nella struttura dei cromosomi}} \\
    \hline
    Cri-du-chat & ritardo mentale, malformazione della laringe \\
    \hline
    \multicolumn{2}{|l|}{\textbf{Cambiamenti nel numero di cromosomi}} \\
    \hline
    Sindrome di Down & ritardo mentale, difetti cardiaci \\
    \hline
    \multicolumn{2}{|l|}{\textbf{Eredità autosomica dominante}} \\
    \hline
    Lobi dell'orecchio liberi & nessuno \\
    \hline
    Acondroplasia & formazione difettiva della cartilagine che causa nanismo \\
    \hline
    Calvizie precoce nei maschi & nessuno \\
    \hline
    Campodattilia & flessione permanente di uno o più dita \\
    \hline
    Capelli ricci & nessuno \\
    \hline
    Malattia di Huntington & progressiva e irreversibile degenerazione del sistema nervoso \\
    \hline
    Sindattilia & fusione di due o più dita \\
    \hline
    Polidattilia & dita sovrannumerarie \\
    \hline
    Brachidattilia & dita più corte del normale \\
    \hline
    Progeria & invecchiamento precoce \\
    \hline
  \end{tabularx}
\end{table}

% ---------------------------------------------------------------------------
\subsection{Poliploidia}

Presenza di un assetto cromosomico multiplo e maggiore di 2 rispetto a quello standard della specie.
\begin{itemize}
  \item deriva da un cattivo funzionamento del fuso durante la mitosi di cellule destinate a diventare germinali;
  \item i cromosomi duplicati restano all'interno di un solo nucleo $\rightarrow$ i gameti prodotti saranno diploidi anziché aploidi;
  \item piante poliploidi: sono più forti e hanno maggior successo nella crescita e nella riproduzione;
  \item animali poliploidi: non vitali; nell'uomo solo l'1\% arriva alla nascita e comunque muore nel corso del primo mese di vita.
\end{itemize}

% ---------------------------------------------------------------------------
\subsection{Mutazioni geniche: mutazioni puntiformi}

Riguardano la sequenza del DNA.

\rubrica{Sostituzione di base}
\begin{itemize}
  \item \textbf{mutazione silente}: la sostituzione di una base non altera l'aminoacido codificato, grazie alla degenerazione del codice genetico;
  \item \textbf{mutazione missenso}: la sostituzione di una base altera l'aminoacido codificato, sostituendolo con un altro;
  \item \textbf{mutazione non senso}: la sostituzione di una base genera un codone di stop, troncando prematuramente la sintesi del polipeptide.
\end{itemize}

\rubrica{Frameshift}
La delezione (o l'inserzione) di uno o più nucleotidi, se non multipli di tre, altera la cornice di lettura a valle del punto della mutazione, cambiando tutti i codoni successivi; può generare un codone di stop prematuro oppure una sequenza aminoacidica completamente diversa da quella normale.

\rubrica{Esempio: anemia falciforme}
La sostituzione di una singola base nel gene della $\beta$-globina determina la sostituzione dell'aminoacido acido glutammico con valina nella catena proteica. Questa mutazione altera la forma dei globuli rossi, che assumono l'aspetto caratteristico ``a falce'' anziché la normale forma discoidale.
